%==============================================================================
%  CAPÍTULO XII
%==============================================================================
\chapter{Intersección del derecho concursal con el derecho laboral y el penal económico}
\label{cap:laboral-penal}

\fichaCapitulo{Efectos en los contratos de trabajo, penales y tributarios.}{Artículo 163 bis
del Código del Trabajo, delitos concursales de la Ley \Nro 21.595, y efectos ante el Servicio
de Impuestos Internos.}

%==============================================================================
\bloqueTeorico
%==============================================================================

\subsection{El derecho laboral y el penal como fronteras del concurso}

El concurso no es un compartimento estanco. Cuando una empresa con trabajadores entra en
insolvencia, la apertura del procedimiento desencadena efectos que el derecho concursal no
regula por sí solo: termina contratos de trabajo, activa preferencias de pago de origen
laboral y, si en la génesis de la insolvencia hubo fraude, abre la puerta a la persecución
penal. El derecho del trabajo y el derecho penal económico operan, así, como las dos
\emph{fronteras} del concurso: los límites donde la lógica concursal cede ante ordenamientos
con principios propios.

Este capítulo aborda esa intersección en su dimensión conceptual y penal. La dimensión
laboral operativa ---terminación de contratos, cálculo y pago de créditos laborales--- se
desarrolla en extenso en los capítulos \ref{cap:laboral-efectos} y
\ref{cap:creditos-laborales}, que integran la Parte dedicada al derecho del trabajo.

\subsection{Fundamento de la preferencia de los créditos laborales}

La preferencia del crédito laboral no es una concesión graciosa: responde a la posición
estructuralmente vulnerable del trabajador frente al empleador insolvente. A diferencia del
acreedor financiero ---que evaluó el riesgo, lo tarificó en la tasa de interés y lo
diversificó en una cartera---, el trabajador no eligió a su deudor, no pudo cubrirse y tiene
en su remuneración un crédito de naturaleza alimentaria del que depende su subsistencia y la
de su familia.

Por eso el ordenamiento sitúa los créditos laborales en la \destacar{primera clase} del
\artcc{2472} \citeyearpar{codigo-civil}, con precedencia sobre los acreedores valistas e
incluso, en ciertos supuestos y con límites, sobre los garantizados con prenda o hipoteca. La
preferencia tiene, con todo, topes ---examinados en el capítulo \ref{cap:creditos-laborales}---
que cumplen una función distributiva dentro de la propia clase, impidiendo que remuneraciones
muy altas absorban el activo en perjuicio de los trabajadores de menor renta.

\begin{foco}[La preferencia laboral condiciona toda la estrategia del concurso]
Antes de proyectar cualquier reparto ---en liquidación--- o cualquier propuesta ---en
reorganización---, debe dimensionarse el pasivo laboral preferente, porque se paga antes que
el resto. Un concurso con fuerte componente laboral deja poco o nada para los valistas, y
tanto el deudor como el acreedor financiero deben saberlo desde el diagnóstico.
\end{foco}

%==============================================================================
\bloqueNormativo
%==============================================================================

\subsection{Área laboral}

\subsubsection{La causal de término del artículo 163 bis}

Causal especial de terminación del contrato de trabajo, que fija como fecha del despido
la de dictación de la resolución de liquidación (\artcdt{163 bis};
\cite{codigo-trabajo}). Sobre el estatuto protector de los trabajadores frente a la
insolvencia del empleador, véase \textcite{doc-proteccion-trabajadores}.

\subsubsection{Preferencia de primera clase del artículo 2472 del Código Civil}

Las remuneraciones, asignaciones familiares y cotizaciones gozan de la preferencia del
\artcc{2472} \Nro 5, que incluye la indemnización del \artcdt{163 bis} \Nro 2 con un tope de
\destacar{90 UF} y los alimentos legales con un tope de \destacar{120 UF} ---valista el exceso---.
Las indemnizaciones legales y convencionales por años de servicio se pagan con la preferencia del
\Nro 8, limitadas a \destacar{tres ingresos mínimos mensuales por año de servicio y fracción
superior a seis meses, con un máximo de once años} (capítulo \ref{cap:creditos-laborales}).

\subsubsection{Finiquitos laborales electrónicos}

Aplicación del instructivo de la \superir{} sobre emisión de finiquitos mediante la
plataforma electrónica de la \dt{} \parencite{instructivo-i6}.\verificar{Número, fecha y vigencia del instructivo.}

\subsection{Área penal: los delitos concursales}

\subsubsection{Ubicación sistemática}

La \LIR{} \citeyearpar{ley-20720} trasladó los delitos concursales desde el antiguo Libro IV
del \ccom{} directamente al \cpen{}, reagrupándolos en el \destacar{Párrafo 7 del Título IX
del Libro Segundo}, que pasó a denominarse «De los delitos concursales y de las
defraudaciones». La \Ldeleco{} \citeyearpar{ley-21595} los incorporó, además, a su catálogo,
con las consecuencias que ese régimen especial acarrea en materia de penas, atenuantes y
determinación.

\subsubsection{Los tipos: artículos 463 y siguientes del Código Penal}

\begin{description}[leftmargin=0pt,style=nextline,font=\sffamily\bfseries]
  \item[Artículo 463 — Reducción del activo] Con pena de \destacar{presidio menor en cualquiera
  de sus grados}, sanciona al deudor que, dentro de los dos años anteriores a la resolución de
  liquidación, conociendo el mal estado de sus negocios, \destacar{disminuye considerablemente su
  patrimonio} ---destruyendo, dañando, inutilizando o dilapidando bienes, o renunciando a créditos
  sin razón--- o ejecuta otro acto manifiestamente contrario a una administración racional. En la
  empresa deudora, la pena se impone también a quien actuó con \destacar{ignorancia inexcusable}
  del mal estado de sus negocios.

  \item[Artículo 463 bis — Ocultación y distracción] Con pena agravada ---\destacar{presidio menor
  en su grado medio a presidio mayor en su grado mínimo}---, sanciona cuatro conductas: favorecer
  a uno o más acreedores pagando deudas no exigibles o constituyendo garantías sobre deudas
  previas dentro de los dos años anteriores; percibir, apropiarse o distraer bienes de la masa
  tras la resolución de liquidación; ejecutar actos de disposición ---reales o simulados--- o
  constituir gravámenes después de esa resolución; y ocultar total o parcialmente sus bienes.

  \item[Artículo 463 ter — Contabilidad e información] Con pena de \destacar{presidio menor en sus
  grados mínimo a medio}, sanciona proporcionar al veedor, liquidador o acreedores información
  falsa o incompleta que no refleje la verdadera situación patrimonial, y la omisión o deficiencia
  en la llevanza y conservación de los libros de contabilidad.

  \item[Artículo 463 quáter — Extensión a administradores] Castiga \destacar{como autor} de los
  delitos anteriores a quien, en la dirección o administración de los negocios del deudor sometido
  a cualquier procedimiento concursal, ejecute los actos o incurra en las omisiones señaladas, o
  los autorice expresamente.
\end{description}

Completan el párrafo el \destacar{artículo 464} ---que sanciona al veedor o liquidador que otorga
ventajas indebidas, con pena de presidio menor máximo a presidio mayor mínimo e inhabilidad
especial perpetua para el cargo--- y las figuras de intervención de terceros (arts. 464 bis y
464 ter). La persecución penal, además, exige \destacar{instancia particular} de la \superir{},
del veedor o liquidador, o de un acreedor que haya verificado su crédito (art. 465): no procede de
oficio.

\begin{foco}[La reforma corrió el eje hacia la administración irracional del patrimonio]
La \Ldeleco{} \citeyearpar{ley-21595} reconfiguró el tipo del artículo 463 del \cpen{}: sanciona hoy al
deudor que ejecuta «actos manifiestamente contrarios a las exigencias de una administración
racional de su patrimonio» que hayan contribuido de manera relevante a ocasionar su
insolvencia. La \destacar{administración irracional del patrimonio} pasa a ser la piedra
angular de la delincuencia concursal, desplazando en parte las hipótesis tradicionales de dolo
e imprudencia. La doctrina de \textcite{mayer-delitos-concursales} analiza en detalle esta
reconfiguración.
\end{foco}


\subsubsection{Los delitos concursales como delitos económicos}

La \Ldeleco{} \citeyearpar{ley-21595} incorporó los ilícitos concursales a su catálogo,
asimilándolos a delitos económicos \destacar{de segunda categoría} cuando se perpetran en el
marco de la dirección o administración de una empresa de mediana o gran envergadura. El régimen
tiene consecuencias severas: \destacar{excluye la aplicación de las atenuantes generales} del
\cpen{} y las sustituye por un sistema propio de atenuantes y agravantes vinculadas al nivel de
culpabilidad y a la existencia de programas de prevención del delito en la organización.

\subsubsection{El requisito de procedibilidad}

Es la nota procesal decisiva: la persecución de estos delitos \destacar{no procede de oficio}.
Sólo puede iniciarse por querella previa de la \superir{}, del veedor o liquidador del
procedimiento respectivo, o de \destacar{cualquier acreedor que hubiere verificado su
crédito}.

\begin{foco}[La verificación abre la puerta penal]
Un acreedor que no verifica su crédito carece de legitimación para querellarse por delito
concursal. La verificación no es sólo un acto de cobro: es el que habilita la acción penal
contra el deudor fraudulento. Para el acreedor que sospecha ocultación, verificar es el
primer paso, no un trámite prescindible.
\end{foco}

\subsubsection{Relación con el incidente de mala fe}

El delito concursal y el incidente de mala fe del \art{169 A} (capítulo \ref{cap:mala-fe})
recaen con frecuencia sobre los mismos hechos ---ocultación, distracción, información falsa---
pero operan en planos distintos: el incidente es sanción civil (denegación del
\discharge{}), con estándar de sana crítica; el delito es sanción penal, con estándar de
certeza. Coexisten sin \emph{non bis in idem}, según se examinó en el capítulo
\ref{cap:mala-fe}. Para la proyección penal del incidente, véase
\textcite{doc-mala-fe-penal}.

\subsubsection{Responsabilidad penal de la persona jurídica}

La \Lrpj{} \citeyearpar{ley-20393} ---cuyo catálogo de delitos base fue notablemente ampliado por
la \Ldeleco{}, que lo abrió a la generalidad de los delitos económicos--- permite perseguir a la
propia empresa cuando los delitos concursales se cometen en su interés o provecho por quienes
ejercen su dirección o administración.

El presupuesto de la imputación es el \destacar{defecto de organización}: la empresa responde
cuando el delito fue posible porque no adoptó los modelos de prevención que le eran exigibles.
Correlativamente, la ley contempla como \destacar{eximente o atenuante} la implementación
efectiva, antes del hecho, de un \destacar{modelo de prevención de delitos} idóneo ---con
encargado de prevención, sistema de supervisión y mecanismos de reporte---.

\begin{foco}[El modelo de prevención como defensa de la empresa]
Para la empresa deudora, el modelo de prevención de delitos no es un formalismo de
\emph{compliance}: es la defensa concreta frente a la imputación penal por delitos concursales de
sus administradores. Una empresa que enfrenta la insolvencia con un modelo de prevención vigente
y operativo tiene un argumento de exención que la que carece de él no puede invocar. La asesoría
concursal de una empresa de tamaño relevante debería revisar el estado de ese modelo desde el
diagnóstico.
\end{foco}

Para una sistematización del régimen tras la reforma, véanse
\textcite{doc-delitos-economicos-garrigues} y \textcite{doc-delitos-economicos-prieto}; sobre la
reconfiguración del tipo base, \textcite{mayer-delitos-concursales}.

\subsection{Área tributaria: efectos del concurso ante el Servicio de Impuestos Internos}

La apertura de la liquidación altera el estatus fiscal del deudor de forma inmediata, y las
consecuencias tributarias son parte del asesoramiento que no puede omitirse.

\subsubsection{Giro inmediato de las diferencias de impuestos}

El artículo 24 inciso cuarto del \ctrib{} faculta al \sii{}, iniciado el concurso,
para \destacar{girar de inmediato} las diferencias de impuestos pendientes de cobro, sin
necesidad de agotar los trámites previos de fiscalización. El crédito fiscal ingresa así al
concurso con celeridad.

\subsubsection{Término de giro simplificado del régimen ProPyme}

La \leyn{21.713}{Ley de Cumplimiento de Obligaciones Tributarias} \citeyearpar{ley-21713}
estableció, para las empresas del régimen ProPyme, un \destacar{procedimiento simplificado y
expedito de término de giro} que reduce los plazos de tramitación administrativa ante el \sii{}.

\subsubsection{Recuperación del IVA impago mediante nota de débito}

El artículo 29 de la \ley{18.591} \citeyearpar{ley-18591} permite a los acreedores
concursales \destacar{recuperar como crédito fiscal el IVA recargado} en facturas impagas
emitidas al deudor en liquidación. Para ello, el liquidador ---en representación de la masa---
debe emitir y notificar la correspondiente \destacar{nota de débito} al acreedor, informando al
\sii{} y a la \tgr{} dentro del período tributario respectivo. Esta vía coexiste con el
mecanismo general del \destacar{artículo 27 ter} del \dl{825} \citeyearpar{dl-825}, que autoriza
la recuperación anticipada del remanente de crédito fiscal originado en documentos incobrables,
cuyo detalle operativo el \sii{} sistematizó en la \textcite{circular-sii-64-2015}. Ambas
herramientas evitan que el IVA recargado y no percibido se transforme en una pérdida definitiva
para el acreedor.

\begin{foco}[El IVA impago no se pierde: se recupera]
El acreedor que emitió factura con IVA a un deudor que luego cayó en liquidación puede
recuperar ese IVA como crédito fiscal, pero necesita la nota de débito del liquidador. Es una
gestión que el liquidador diligente ofrece y que el acreedor debe reclamar: sin la nota de
débito, el beneficio del artículo 29 de la Ley 18.591 no opera.
\end{foco}

\subsubsection{Castigo de créditos incobrables y tratamiento del saldo insoluto}

Para el acreedor que no alcanza reparto, el crédito verificado y no pagado en la liquidación se
torna \destacar{incobrable} y puede castigarse como gasto conforme al \artir{31}{4}, con las
exigencias de agotamiento de cobro que el \sii{} precisó en la \textcite{circular-sii-62-2016}.
La resolución de liquidación y el certificado de incobrabilidad del liquidador operan aquí como
antecedentes de respaldo del castigo.

Correlativamente, para la persona deudora, la \destacar{extinción de los saldos insolutos} que
produce el \discharge{} de la liquidación no constituye un incremento de patrimonio gravable:
la remisión legal de deudas queda comprendida entre los ingresos no constitutivos de renta del
\artir{17}{22}. El deudor rehabilitado no tributa por la deuda que la ley extinguió.

\begin{practica}[La preferencia fiscal del art. 2472 N.\texorpdfstring{\textsuperscript{o}}{o} 9]
No todo crédito del Fisco es valista. Los impuestos de \destacar{retención y recargo} ---IVA,
retenciones del impuesto único de segunda categoría--- gozan de la preferencia de primera clase
del \artcc{2472} \Nro 9, con precedencia sobre los acreedores comunes. Distinguir el tramo
preferente del tramo valista del crédito fiscal es decisivo al proyectar el reparto.
\end{practica}

%==============================================================================
\bloquePractico
%==============================================================================

\subsection{Desvinculación colectiva de los trabajadores}

\begin{flujograma}{Secuencia de desvinculación colectiva tras la resolución de liquidación}{cap12-despidos}
  \node[hito] (sent) {SENTENCIA DE LIQUIDACIÓN DE LA EMPRESA\\[2pt]
    \normalfont Se configura la causal del \artcdt{163 bis}};
  \node[nodo, below=of sent] (comunica)
    {EL LIQUIDADOR COMUNICA EL DESPIDO A LOS TRABAJADORES Y A LA INSPECCIÓN DEL TRABAJO\\[2pt]
     \normalfont Plazo legal de 6 días desde la dictación};
  \node[nodo, below=of comunica] (liquida)
    {LIQUIDACIÓN INDIVIDUAL DE CRÉDITOS Y PREFERENCIAS LABORALES};
  \node[nodo, below=of liquida] (finiquito)
    {EMISIÓN Y SUSCRIPCIÓN DEL FINIQUITO ELECTRÓNICO\\[2pt]
     \normalfont Portal de la \dt{}};
  \node[favorable, below=of finiquito] (pago)
    {PAGO CON PREFERENCIA DE PRIMERA CLASE CON RECURSOS DE LA MASA};

  \draw[flecha] (sent) -- (comunica);
  \draw[flecha] (comunica) -- (liquida);
  \draw[flecha] (liquida) -- (finiquito);
  \draw[flecha] (finiquito) -- (pago);
\end{flujograma}

\begin{alerta}[Plazo de pago del finiquito]
El plazo que la \dt{} aplica ordinariamente para el pago de los finiquitos suscritos en
la plataforma no resulta vinculante en sede concursal: el cobro de las indemnizaciones se
sujeta al estado de realización de los activos de la masa.\verificar{Confirmar que el
instructivo vigente mantiene esta regla.}
\end{alerta}

%==============================================================================
\bloqueJurisprudencial
%==============================================================================

\subsection{Fuero laboral y despido indirecto en sede concursal}

Criterios de los tribunales laborales sobre la vigencia del fuero maternal y sindical de
trabajadores de empresas declaradas en liquidación, y el efecto terminador del
\artcdt{163 bis}.\pendiente{Sistematizar jurisprudencia de unificación.}

\subsection{Responsabilidad de gerentes y directores}

Jurisprudencia que delimita la responsabilidad penal por ocultación o distracción de
activos en los meses previos a la apertura del concurso.\pendiente{Sistematizar.}
